\subsection{Finite local contradictions and their certificates}
\label{app:finite-certificate}
The halting switch uses a fixed finite computation prefix even though
function totality is expressed by sums with $n$ summands. The following
elementary certificate construction separates these two sizes.

\begin{lemma}[Finite-term certificate]\label{lem:finite-certificate}
Let $n\ge1$ and let a finite acyclic list of $m$ function evaluations be given, with every
argument supplied by an earlier node. Fix $b$ distinct labels. Suppose that
every assignment of labels to the nodes, and every Boolean assignment to
the relation atoms on the named and fixed labels, either has inconsistent
function values or falsifies a clause instance whose variables lie among the named and fixed
labels and all of whose function evaluations occur in the list. If clauses have length at most $L$, the
graph-and-sum translation has an NS refutation of degree at most
\[
                 m+\sum_R(m+b)^{\operatorname{arity}(R)}+L.
\]
The hypothesis and conclusion concern sizes for which the $b$ distinct
labels exist.
\end{lemma}
\begin{proof}
For a labelling $\alpha$ of the nodes, let $q_\alpha$ be the product of the
$m$ graph variables asserting those function evaluations. Expanding the
totality sum for each node, in the order of the acyclic list, gives
\[
                         \sum_\alpha q_\alpha-1
\]
in the degree-$m$ NS span. This is a telescoping identity: at step $j$, the
product for the first $j-1$ evaluations multiplies the totality equation
for evaluation $j$. Summing over the possible earlier values gives the
next partial sum. Repeated graph variables cause no problem.

If $\alpha$ gives different values to the same function on the same input,
$q_\alpha$ contains a functionality collision and already belongs to the
span within degree $m$. Otherwise, list the distinct relation variables
on the at most $m+b$ named labels. Write this actual set as
$Y_\alpha=\{y_1,\ldots,y_{a_\alpha}\}$, where
$a_\alpha\leq a_{\mathrm{rel}}=\sum_R(m+b)^{\operatorname{arity}(R)}$.
Expand over these distinct variables:
\[
             1=\prod_{i=1}^{a_\alpha}(y_i+(1-y_i))
              =\sum_\beta u_\beta.
\]
Each $u_\beta$ is the Boolean assignment monomial for one partial relational
structure. By hypothesis that structure falsifies an available clause.
Its clause polynomial is a product of at most $L$ falsity indicators.
Write $P=q_\alpha u_\beta$ and the clause polynomial as
$F_1\cdots F_\ell$, where each $F_j$ is the falsity indicator of a
literal false in the chosen partial structure. For each literal,
$P(1-F_j)$ has an explicit expression in the structural NS span:
\begin{itemize}
\item For a positive relation literal, $F_j=1-x$ and $u_\beta=(1-x)u'$.
Thus $P(1-F_j)=q_\alpha u'x(1-x)=-q_\alpha u'(x^2-x)$.
For a negative relation literal, $F_j=x$ and $u_\beta=xu'$; the same
identity with $x(1-x)$ applies.
\item For a negative graph literal false in the partial structure,
its graph variable $x$ is one of the chosen evaluations, so
$q_\alpha=xq'$. Its falsity indicator is $F_j=x$ and
$P(1-F_j)=-q'u_\beta(x^2-x)$.
\item For a positive graph literal false there, the chosen evaluation
has a different output variable $x'$ on the same function input.
Here $q_\alpha=x'q'$, $F_j=1-x$, and
$P(1-F_j)=q'u_\beta x'x$, an explicit functionality-axiom multiple.
\item A false equality or disequality has the constant falsity
indicator $F_j=1$, so $P(1-F_j)=0$.
\end{itemize}
These are all cases of the normalized literals. Each identity has
product degree at most $m+a_{\mathrm{rel}}+1$.
The exact telescoping identity
\[
 P(1-F_1\cdots F_\ell)
   =\sum_{j=1}^{\ell}P(1-F_j)\prod_{i<j}F_i
\]
therefore puts its left side in the NS span within degree
$m+a_{\mathrm{rel}}+\ell$. Since $PF_1\cdots F_\ell$ is a
multiple of the falsified clause axiom, $P$ is in that same span.
Sum first over $\beta$, then over $\alpha$, and use the totality identity
to obtain one.
\end{proof}


A \emph{partial structure} for $\theta_M$ has a domain of labels, all relations defined on it, and
function values defined only on some arguments. It is \emph{locally consistent} if every clause
instance whose variables take values in the domain, and whose function atoms are all defined,
has a true literal.

\begin{lemma}[Local switch]\label{lem:local-switch}
Let the run of $M$ first reach a halting state at time $T$, and let $M$ have $\kappa$ counters.
Suppose a locally consistent partial structure defines $o$, $e$, $a_{j+1}=S(a_j)$ for $j<T$
\textup{(}with $a_0=o$\textup{)}, $C_r(a_j)$ and $S(C_r(a_j))$ for $j\le T$, $r<\kappa$, and that all these
elements lie in its domain. Then its domain has at most $T$ distinct elements.
\end{lemma}
\begin{proof}
Suppose it contains $T+1$ distinct elements. If $a_j=e$ for some $j<T$, take $j$ least. Every
domain element $x$ satisfies $x=a_0$ or $a_0<x$ (least element), and, if $a_i<x$ with $i<j$, then
$x=a_{i+1}$ or $a_{i+1}<x$ (no element strictly between $a_i$ and $S(a_i)$, and linearity). As
$a_j=e$ is greatest, the domain is $\{a_0,\ldots,a_j\}$, too small. So $a_j\ne e$ for $j<T$, and
the order clauses make $a_0<a_1<\cdots<a_T$. For the counter induction, the initial clauses give the initial
state and $C_r(a_0)=a_0$ for every counter. A counter value at time $j$
is at most $j$, so the notation $a_v$ uses only this named prefix.
Suppose its states and counters agree through time $j<T$.
An increment of value $v$ gives $S(a_v)=a_{v+1}$; because $v\le j<T$,
this successor is among the defined evaluations and its guard is valid.
The zero branch of a decrement is selected exactly when $v=0$.
On a positive decrement, the transition specifies a value $y<a_v$
with $S(y)=a_v$. The immediate-predecessor property forces
$y=a_{v-1}$: if $y<a_{v-1}$ then $a_{v-1}$ lies strictly between
$y$ and $S(y)$, and if $a_{v-1}<y<a_v$ then $y$ lies strictly
between $a_{v-1}$ and $S(a_{v-1})$. Both alternatives contradict
immediacy. Unchanged counters retain their named values. State
uniqueness and the chosen transition identify the next state.
Every clause used has its variables in the domain and its function
atoms in the specified finite evaluation list. This proves the
induction through time $T$. At $j=T$ the state is halting, which a clause forbids.
\end{proof}

\begin{corollary}\label{cor:switch-degree}
In that case, for every $n\ge T+1$ the translation of $\theta_M$ at size $n$ has an NS refutation
of degree
\[
 C_M=m+\sum_R(m+T+1)^{\operatorname{arity}(R)}+\max_{C}|C|,\qquad m=T+2+2\kappa(T+1).
\]
\end{corollary}
\begin{proof}
Apply Lemma~\ref{lem:finite-certificate} to the $m$ evaluations $o$, $e$, $a_1,\ldots,a_T$,
$C_r(a_j)$, $S(C_r(a_j))$ ($j\le T$, $r<\kappa$), each argument supplied by an earlier node, with
$b=T+1$ fixed labels. A consistent labelling and a Boolean assignment give a partial structure of
the kind in Lemma~\ref{lem:local-switch}, whose domain contains the $T+1$ fixed labels; so it is not
locally consistent, which is the hypothesis of Lemma~\ref{lem:finite-certificate}.
\end{proof}

Here $\kappa$ is one more than the largest counter index occurring in $M$ (so $\kappa\ge1$), and $T$ must be
the first halting time. Sizes $n\le T$ are genuine exceptions: $\theta_M$ can have models there.
